\documentclass[11pt]{article}
\usepackage[margin=1in]{geometry}
\usepackage{amsmath,amssymb,mathtools}
\usepackage{enumitem}
\usepackage{microtype}

\title{Two Neural Approaches to the Stationary Kyle--Back Lag System}
\author{}
\date{\today}

\newcommand{\Dt}{\widetilde{\mathcal D}}
\newcommand{\Vt}{\widetilde V}
\newcommand{\calD}{\mathcal D}
\newcommand{\dd}{\,d}

\begin{document}
\maketitle

\section*{Purpose}

There are two numerically different ways to use a neural network in the
stationary Kyle--Back system.  They are easy to confuse because both are
``PINN''-like, both work on stationary lags, and both use the same direct
observation geometry.  The difference is where the transport and birth
structure is enforced.

The current full PINN asks the network to learn the observer-specific unresolved
kernels
\[
  \Vt^a(\ell), \qquad \Dt^{\mathrm{tot},a}(\ell),
\]
the policy/readout kernels, and sometimes adjoints, then penalizes the filter,
projection, and first-order-condition residuals.  The lag-reduced approach fixes
the candidate policy/filter environment and learns only a few scalar residual
readouts.  It reconstructs the two-lag noise-state response analytically from
the stationary transport equation and its diagonal birth traces.

\section*{Common observation geometry}

Use primitive Brownian coordinates
\[
  W=(W^V,W^Z,W^{Y1},W^{Y2}).
\]
The market maker directly observes the public order-flow noise subspace
\[
  \Pi^0=E^{Z,\top}E^Z,
\]
while trader \(i\) directly observes order-flow noise plus their own private
signal noise,
\[
  \Pi^1=E^{Z,\top}E^Z+E^{Y1,\top}E^{Y1},\qquad
  \Pi^2=E^{Z,\top}E^Z+E^{Y2,\top}E^{Y2}.
\]
There is no direct birth term for the fundamental-value Brownian coordinate
\(W^V\).  At zero lag, indirect old-history terms vanish; only the appropriate
directly observed innovation row can appear through the birth trace.

\section*{Approach A: full unresolved-kernel PINN}

The full solver represents the unresolved kernels directly.  For observer
\(a\), the order-flow and value residuals are written as
\[
  \Dt^{\mathrm{tot},a}_t(u)
  =
  \calD^{\mathrm{tot}}_t(u)(I-\Pi^a)
  -
  \int_{-\infty}^t
    \calD^{\mathrm{tot}}_t(r)F^a_t(r,u)\dd r,
\]
\[
  \Vt^a_t(u)
  =
  \Sigma_V^{1/2}E^V(I-\Pi^a)
  -
  \int_{-\infty}^t
    \Sigma_V^{1/2}E^V F^a_t(r,u)\dd r .
\]
The normalized unresolved observation rows are then
\[
  \widetilde C^0(u)=\Sigma_Z^{-1/2}\Dt^{\mathrm{tot},0}(u),
\]
\[
  \widetilde C^i(u)=
  \begin{pmatrix}
    \Sigma_Z^{-1/2}\Dt^{\mathrm{tot},i}(u)\\
    \Sigma_{Y,i}^{-1/2}\Gamma_i\Vt^i(u)
  \end{pmatrix}.
\]
The filter kernel \(F^a\) is built from these rows by the characteristic
formula with side births:
\[
  F^a(\ell,m)
  =
  \mathbf 1_{\ell>m}\,\widetilde C^a(\ell-m)^\top E^a
  +
  \mathbf 1_{m>\ell}\,E^{a,\top}\widetilde C^a(m-\ell)
  +
  \int_0^{\min(\ell,m)}
    \widetilde C^a(\ell-\tau)^\top
    \widetilde C^a(m-\tau)\dd\tau .
\]
The network is penalized for failing the closure identities above, for failing
projection identities such as \(P^a K=K\Pi^a+\int K F^a\), and, in the coupled
version, for failing the weak stationary first-order condition.

This formulation is closest to the chapter equations.  Its weakness is that a
large neural network can satisfy collocation points while producing high
frequency oscillations in weakly identified rows.  In particular, the
transport/birth geometry is represented only through residual penalties, so
endpoint leakage or insufficient randomization can produce spurious wiggles.

\section*{Approach B: lag-reduced forward IRF PINN}

The lag-reduced approach starts from a fixed candidate environment:
policy kernels \(D^1,D^2\) and unresolved observation rows
\[
  c^0,\qquad c^{j,Z},\qquad c^{j,Y}.
\]
For a unit deviation/spike by trader \(i\), let \(j\neq i\).  The network does
not learn a full two-lag kernel.  It learns only three scalar residual readouts
as functions of future lag \(s\):
\[
  r_i(s)=\bigl(r_0(s),r_Z(s),r_Y(s)\bigr).
\]
These are the coefficients feeding the market-maker order-flow residual row,
the opponent's order-flow row, and the opponent's private-signal row.

Given \(r_i\), the old-history density for future lag \(s\) and age \(a\) is
constructed explicitly.  Schematically,
\[
\begin{aligned}
U_0(s,a)
={}&
\mathbf 1_{a\ge s}\,\frac{c^0(a-s)}{\sigma_Z}
+
\mathbf 1_{a<s}\,r_0(s-a)E^Z\\
&+
\int_{\max(0,s-a)}^s
  r_0(u)c^0(a-s+u)\dd u ,
\end{aligned}
\]
and for the opponent observer,
\[
\begin{aligned}
U_j(s,a)
={}&
\mathbf 1_{a\ge s}\,\frac{c^{j,Z}(a-s)}{\sigma_Z}
+
\mathbf 1_{a<s}
  \{r_Z(s-a)E^Z+r_Y(s-a)\sigma_{Y,j}E^{Yj}\}\\
&+
\int_{\max(0,s-a)}^s
  \{r_Z(u)c^{j,Z}(a-s+u)
    +r_Y(u)c^{j,Y}(a-s+u)\}\dd u .
\end{aligned}
\]
The diagonal atom is also reconstructed from the same residuals, for example
\[
 A_0(s)
 =
 \frac{E^Z}{\sigma_Z}
 +
 \int_0^s r_0(u)c^0(s-u)\dd u,
\]
\[
 A_j(s)
 =
 \frac{E^Z}{\sigma_Z}
 +
 \int_0^s
 \{r_Z(u)c^{j,Z}(s-u)+r_Y(u)c^{j,Y}(s-u)\}\dd u .
\]
The zero-lag structure is therefore hard coded: there is no value-noise birth,
trader 1 never receives a \(Y2\) birth, and trader 2 never receives a \(Y1\)
birth.

The loss then asks the learned residual readout to be self-consistent.  If
\(\Delta D_j(s)\) is the induced opponent-demand response,
\[
  \Delta D_j(s)
  =
  \int U_j(s,a)\cdot D^j(a)\dd a
  +
  A_j(s)\cdot D^j(s),
\]
and the self projections from the market-maker and opponent observation rows
are
\[
  S_0(s)=
  \int U_0(s,a)\cdot c^0(a)\dd a + A_0(s)\cdot c^0(s),
\]
\[
  S_Z(s)=
  \int U_j(s,a)\cdot c^{j,Z}(a)\dd a + A_j(s)\cdot c^{j,Z}(s),
\]
\[
  S_Y(s)=
  \int U_j(s,a)\cdot c^{j,Y}(a)\dd a + A_j(s)\cdot c^{j,Y}(s),
\]
then the target residual is
\[
  r_0(s)=\frac{\Delta D_j(s)}{\sigma_Z}-S_0(s),\qquad
  r_Z(s)=\frac{\Delta D_j(s)}{\sigma_Z}-S_Z(s),\qquad
  r_Y(s)=-S_Y(s).
\]

\section*{Practical consequences}

\begin{itemize}[leftmargin=*]
\item The full unresolved-kernel PINN is a candidate equilibrium solver.  It can
move filters, policies, and adjoints together, but it is high-dimensional and
more vulnerable to oscillatory weak solutions.
\item The lag-reduced PINN is a forward IRF solver for a fixed environment.  It
is lower-dimensional and enforces the transport and birth equations by
construction, so it should be much less prone to artificial two-lag
oscillation.
\item The lag-reduced solver is not, by itself, a complete equilibrium solver.
To turn it into one, it should be embedded in an outer loop that updates the
candidate policy/filter environment and retrains or resolves the residual
readouts.
\item For the current debugging problem, the lag-reduced version is the cleaner
diagnostic.  If it still oscillates, the issue is probably in the candidate
equations or readout target.  If it does not oscillate while the full PINN does,
the issue is likely representation/identification in the full unresolved-kernel
network.
\end{itemize}

\end{document}
